\chapter{Search}

``Nani, can you please give me eleven rupees? I need to recharge my phone,'' Shiva asked, sitting in the verandah.

``Didn't Nana get it recharged just last week?'' Nani replied, carrying on with the knitting she had probably been at for the past month.

``Yes, but I wanted an internet recharge. It's separate from the calling one—it's for searching things on the internet,'' he said, sliding closer to her on the sofa.

``What is internet?'' she asked innocently, her eyes still on the needles.

``It's like a whole world inside the phone that holds all the information about everything there is. You write down your questions and it answers them.''

``Oh god, the things people have made.'' She stopped knitting and turned to him. ``Can it tell me where my earrings are? I kept them in one of the boxes and now I can't find them.''

``No, Nani,'' Shiva replied with a sigh, ``it only knows what people have put on it for others to see.''

``Ohh.'' Her enthusiasm drained away, and she went back to her knitting. ``Then what's the use? You could just ask the person who put it there in the first place.''

``But how would I know who knows the answer?'' Shiva said, confident, as if the argument alone could win it for him.

``Shhoo... shooo...'' Nani called out sharply at a cat trying to slip inside with its kitten. Shiva got up and moved quickly toward the gate; at the sight of him, the cat darted out.

``All right—go to the almirah and take it from there,'' Nani said, returning to her knitting.

``Sure. I know where it is.''

``And once you get your internet,'' she added with a smile, ``show it to me too—what it looks like.''

Shiva hurried to the nearby shop, bought a recharge coupon, and came back home, his steps quick—as though something were slipping away and he had to reach it before it was gone. He came in through the gate, Nani still sitting where he'd left her, and rushed straight past her.

``Where are you going? Did you get it?'' she asked.

Shiva stopped just behind her. ``Yes. I'm going upstairs—the network stays weak down in the house.''

``But it'll be cold up there, look how breezy it is,'' Nani said, concern in her voice. ``You'll fall sick. Why not go up some other day, when there's a bit of sun?''

``No, it's fine, Nani. I'll stay in behind the bulkhead.''

``Still—wear your jacket on top,'' she added.

It was truly windy, just as Nani had said. Down on the street the wind was weak, broken up by the houses all around, but three floors up it had nothing to stop it. Not that anything could have stopped Shiva. He scratched the voucher and got the 20 MB of data he'd been so impatient for, opened GPRS, went to Facebook, and made himself a profile—just his name on it. A profile picture? Forget it; the phone couldn't even render the Facebook logo.

He opened the search bar and typed ``Gauri Mishra.'' A long list of girls came up, not one of them with a picture. Now that was a problem—how was he to know which one was really her? He read through every bio; none carried any mark of his school, or even his town. He was stuck. Then a thought struck him, and he began sending ``Hi'' to all of them, one after another—two hours of patience in the cold, until his phone finally died.

``You know, I tried to reach you so many times all evening, but your phone was off. What were you even doing? You know how little time I get to talk to you,'' Shreya said, her voice sharp with anger.

``I'm sorry. The phone just died—I put it to charge and forgot to switch it back on. I didn't mean to hurt you,'' he said, defending himself.

``It's not about that. How can you forget that I call you in the evening?'' she asked, her voice rising.

``I'm really sorry—but please, keep your voice down. It's one in the morning; someone could wake up and hear you,'' he said, worried.

``I know. That's why I'm on the roof. No one can hear me up here.''

``What? Why are you on the roof at this hour? It's freezing outside. Please come down—we can talk tomorrow,'' he said, his voice growing more concerned.

``It's fine. Bother me as much as you like while you still have me. Because if I'm ever gone, you won't find anyone who loves you the way I do.'' Her voice grew sadder, close to tears.

``Hey—why are you crying? I'm right here. I'm not going anywhere. Don't worry.''

``I'm not crying, it's just my nose running,'' she said, clearly wiping her tears and trying to stop sobbing.

``Please, don't hide from me. I know you. I'm really sorry—I didn't mean to hurt you. Please forgive me,'' he said, as gently as he could.

``It's okay. Just remember to answer when I call. I really don't get much time with you.''

``Understood, ma'am. You command my life,'' he said, joking.

``Now please, come down and get some sleep. You'll miss school tomorrow,'' he went on.

``No, I won't. Someone will wake me in the morning.''

``Oh—I forgot, my princess has her servants too.'' They both cracked up.

``So—have you confirmed you're coming for the pooja? It's less than a month away now,'' he asked.

``Yes, sir, I will. And I'm going shopping this weekend—tell me your favourite colour and I'll buy something for myself in it.''

``It's black. But I don't think it would suit you. Take something in crimson—it would look best on you.''

``Sure, I'll see,'' she replied, with a final wipe of her tears. ``Okay, I have to go now, or I'll get dark circles—and I'm telling you, I don't look good with those,'' she added, smiling.

``Sure, sure...'' Shiva said, and they both laughed before hanging up.

``Wake up. Shiva, wake up. We have to leave early.''

Shiva opened his eyes to find Mama standing beside his bed. Rubbing them, he asked, ``What happened?''

``Let's get the water pump out to the field. I've found someone to help me through the day, and you can leave for school once we've set it up,'' Mama said, taking a sip of tea from his cup.

Shiva turned to his phone and checked the time. ``But it's five in the morning.''

``Yes—Shirilal will be here by five thirty. I've already carried the pump set out to the gate. Get ready.''

``Okay.'' Shiva pulled himself out of the warm bed into the cold and headed for the bathroom.

He came back to the verandah a few minutes later.

``Babu, I've left the tea on the table with some snacks—eat before you go,'' Mami said, heading back inside through the gallery.

Shiva picked up the tray and went out to the porch, where Mama stood holding his cup of tea, inspecting the pump set.

``Isn't Shirilal the one who joined the police? A constable?'' Shiva asked, taking his first sip from the cup. ``Rohan said he sat the medical exam and had cleared every stage before it.''

``He did,'' Mama said, setting his empty cup aside and crouching by the machine to inspect something only he understood. ``But his family couldn't put together the ten lakhs the recruiters wanted. So he had to turn back.''

``But he's good at his studies, and his fitness too—I heard he's one of the fastest runners of everyone who trains at the Tedhiya orchard.'' Shiva picked a pakora off the tray.

``Talent isn't the whole of it, Shiva.'' Mama said with a sigh, already fixing something with a wrench.

``Namaste, uncle,'' a voice called from across the street. Shiva looked round to see Shirilal coming toward them.

``Have some tea, Shirilal, and then we'll go,'' Mama said without turning.

``No, uncle, it's fine. I had breakfast at home—my mother made it for me,'' he replied, his voice bright, not a trace of the loss he'd been through.

``Shiva, finish up quickly,'' Mama said, wiping the oil off the part he'd just fixed.

Soon Shiva and Shirilal took up the pump set, Mama wheeled the bicycle with the pipes loaded on it, and they set off for the field. Dawn was breaking on the horizon, a reddish hue glazing the riverbed and the mustard field. Across the river, the far village seemed dim beyond the dazzle of the water, the people moving about it small as stick figures—almost cartoonish as they crossed the fields.

``Shirilal, Shiva and I will get the pump set up. You take the bucket and fetch water for the suction,'' Mama said, and began untying the pipes. Shiva joined in, dragging the pump toward the bore.

``But Mama—what will he do now? He can't just go on doing manual labour after finishing his graduation,'' Shiva asked, his eyes full of concern.

``He's still looking for other government jobs.''

``And if the same thing happens again?''

``Can we stop trying just for fear of failing?'' Mama said, fitting the pipe to the pump and the bore. ``What if we fail? We keep searching for what we want. Who knows what God has kept hidden for us.''

Shiva nodded, watching Shirilal come toward them along the narrow farm trails, a bucket in each hand.